\section{Q2 further results}\label{sec:app-q2}
\subsection{Wallet and volume distributions}\label{sec:app-q2-dist}
% src: S-PAPER-5, S-PAPER-10
\begin{figure}[H]
\centering
\begin{subfigure}[t]{0.49\linewidth}\centering
\figorbox{figures/wallet_dist_maker_count}{\linewidth}
\caption{Share of maker wallets per decade of volume.}
\label{fig:wd-maker-count}
\end{subfigure}\hfill
\begin{subfigure}[t]{0.49\linewidth}\centering
\figorbox{figures/wallet_dist_maker_volume}{\linewidth}
\caption{Share of maker volume per decade of wallet volume.}
\label{fig:wd-maker-volume}
\end{subfigure}
\caption{Wallet and volume distributions of makers, split at the top \TopTwentyLabel{} of wallets. Why: the area is the share of wallets in (a) and the share of volume in (b).}% Source: run \WalletDistRun.
\label{fig:wd-maker}
\end{figure}

\begin{figure}[H]
\centering
\begin{subfigure}[t]{0.49\linewidth}\centering
\figorbox{figures/wallet_dist_taker_count}{\linewidth}
\caption{Share of taker wallets per decade of volume.}
\label{fig:wd-taker-count}
\end{subfigure}\hfill
\begin{subfigure}[t]{0.49\linewidth}\centering
\figorbox{figures/wallet_dist_taker_volume}{\linewidth}
\caption{Share of taker volume per decade of wallet volume.}
\label{fig:wd-taker-volume}
\end{subfigure}
\caption{Wallet and volume distributions of takers, split at the top \TopTwentyLabel{} of wallets. Why: the area is the share of wallets in (a) and the share of volume in (b).}% Source: run \WalletDistRun.
\label{fig:wd-taker}
\end{figure}

\begin{figure}[H]
\centering
\begin{subfigure}[t]{0.49\linewidth}\centering
\figorbox{figures/wallet_lorenz_maker}{\linewidth}
\caption{Makers.}
\label{fig:lorenz-maker}
\end{subfigure}\hfill
\begin{subfigure}[t]{0.49\linewidth}\centering
\figorbox{figures/wallet_lorenz_taker}{\linewidth}
\caption{Takers.}
\label{fig:lorenz-taker}
\end{subfigure}
\caption{Lorenz curves of volume, with the 80/20 point. Why: the height is the volume share of the lowest wallets. }% Source: run \WalletDistRun.
\label{fig:lorenz}
\end{figure}

\subsection{k-means side analysis}\label{sec:app-kmeans}
% src: S-PAPER-2, S-PAPER-6
We take the logarithm of the heavy-tailed features, then z-score them.
We choose $K$ by a fixed rule on the elbow curve.
The rule takes the point farthest from the line between the first and last points.
We test cluster against topic category with the chi-square test of independence .
Its effect size is the largest observed minus expected share in one cell.
This analysis is exploratory.

% src: S-PAPER-5, S-PAPER-10
In the side analysis, the fixed elbow-curve rule picks $K=\KmeansK$.
The sample has \KmeansNSample{} wallet-role rows, and a wallet in both roles gives two rows.
The \KmeansK{} clusters have \ClusterZeroN, \ClusterOneN, \ClusterTwoN, and \ClusterThreeN{} rows.
Their median volumes run from \ClusterTwoMedianVolume{} to \ClusterOneMedianVolume{} USDC.
An ANOVA of mean trade price across clusters gives $\mathrm{SSB}/(\mathrm{SSB}+\mathrm{SSW})=\KmeansPriceSsbRatio$.
Its permutation $p$ is at the floor, $1/(\KmeansNPerm+1)$.
The cluster by category test gives $\chi^2=\KmeansCategoryChi$ with \KmeansCategoryDf{} degrees of freedom.
Its $p\,\KmeansCategoryP$ is a float-underflow bound, on \KmeansNSample{} rows.
The largest cell gap is in cluster \ClusterCategoryGapCluster, category \ClusterCategoryGapCategory.
Its observed share is \ClusterCategoryGapObs{} of the sample, against \ClusterCategoryGapExp{} expected.
The gap is \ClusterCategoryGapDiff.

\subsection{Wallet profit and loss}\label{sec:app-q2-pnl}
% src: S-PAPER-5, S-PAPER-10
This analysis is descriptive and outside the pre-registered Q2 criteria.
It asks how much each wallet won or lost on resolved markets.
It uses the kept markets only.
The gross profit of a wallet is its cash flow plus its payout, summed over all its trades in the scope.
Its net profit is the gross profit minus the recorded taker fee of each trade.
A buy costs the USDC amount and gives shares at the trade price.
A sell gives the USDC amount and takes shares.
The maker takes the opposite leg of the same trade.
Each share pays 1 USDC if its outcome won and 0 otherwise.
A wallet is a winner if its net profit is above $10^{-6}$ USDC, a loser if below $-10^{-6}$ USDC, and zero otherwise.
We report no fee totals, because the meaning of the fee column is unverified.
The net profit rests on that column.

\begin{figure}[H]
\centering
\figorbox{figures/fig_pnl_distribution}{0.62\linewidth}
\caption{Wallets by net profit in symmetric-log bins, for all resolved markets, for markets without Up/Down, and for the kept markets (bottom panel). Why: it shows how many wallets sit near zero and how long both tails are.}% Source: run \PnlRun.
\label{fig:pnl-dist}
\end{figure}

% src: S-PAPER-5, S-PAPER-10
In the kept markets, \PnlNWallets{} wallets include \PnlNWinners{} winners and \PnlNLosers{} losers.
Wallets won \PnlGrossWonM{} M USDC and lost \PnlGrossLostM{} M USDC before fees.
The per-market identity check, that gross profits over all wallets sum to zero, had \PnlNIdentityFailures{} failures.
It covers \PnlNMarketsChecked{} usable markets in the full scope, not only the kept markets.
The median net profit is \PnlMedianNet{} USD.
The 1\% quantile is \PnlQuantileOne{} USD and the 99\% quantile is \PnlQuantileNinetyNine{} USD.
The top 1\% of wallets take an estimated \PnlTopOnePctShareWon{} of all winnings.
We use quantiles and a histogram, which are in the lectures.

The wallet-level numbers are estimates, for these reasons.
\begin{itemize}
\item A trade matched by a mint or a merge has no wallet on the other side. We do not model it, and \PnlShareShortPositions{} of kept positions are short.
\item \PnlNKeptNotUsable{} kept markets fail the winner rule, so their payouts are not counted.
\item Every position is held to resolution, and no wallet sells after resolution.
\item Duplicate rows stay, so a wallet's profit includes the repeats.
\item Fees are not reported.
\end{itemize}

